\documentclass[11pt]{article}

\usepackage[a4paper,margin=1in]{geometry}
\usepackage{amsmath,amssymb}
\usepackage[hidelinks]{hyperref}
\usepackage{enumitem}

\title{Second Codex Review Findings for\\
\texttt{rational\_cut\_and\_project\_gap\_periods.tex}}
\author{Codex Review Notes}
\date{May 7, 2026}

\begin{document}

\maketitle

\section*{Executive Summary}

I reviewed the current manuscript
\[
\texttt{LaTeX/rational\_cut\_and\_project\_gap\_periods.tex}
\]
again after the revisions based on the previous findings.  The substantial
issues from the first review appear to have been addressed.  In particular:

\begin{itemize}[leftmargin=*]
  \item the residue-bijection corollary now has a proof;
  \item the bridge from gap-periodicity to residue-multiplicity invariance is
        now explicit;
  \item the CSV verification data now match both the multiset and set formulas;
  \item the verification script now handles the current six-column CSV format;
  \item the stale Lean build-log files have been removed;
  \item the paper compiles cleanly with \texttt{latexmk};
  \item the Lean project builds successfully.
\end{itemize}

I did not find a new mathematical objection to the rational main theorem or to
the set-valued companion theorem.  The remaining points are presentation and
reproducibility polish before arXiv submission.

\section*{Remaining Findings}

\begin{enumerate}[label=\textbf{F\arabic*.},leftmargin=*]

\item \textbf{The Lean table has a stale commit reference.}

In the Lean correspondence table caption, around
\[
\texttt{LaTeX/rational\_cut\_and\_project\_gap\_periods.tex:1465--1468},
\]
the manuscript says that the line numbers refer to commit
\[
\texttt{0b12c26}.
\]
The current repository head at the time of this review is
\[
\texttt{e8cb69b}.
\]
Moreover, at commit \texttt{0b12c26} the Lean directory was still named
\[
\texttt{LEAN/CutAndProject/...},
\]
whereas the manuscript now correctly refers to
\[
\texttt{Lean/CutAndProject/CutAndProject/Basic.lean}.
\]

\textbf{Recommendation.}  Either update the commit hash to the current
publication commit, or remove the commit hash from the table caption and say
that the line numbers refer to the companion repository version submitted
with the paper.

\item \textbf{Remove internal project language from the arXiv text.}

Around
\[
\texttt{LaTeX/rational\_cut\_and\_project\_gap\_periods.tex:1544--1551},
\]
the paper says:
\[
\text{``After Phase~2 of the multiplier cleanup, ...''}
\]
This reads like an internal development note rather than publication prose.

\textbf{Recommendation.}  Replace it by a neutral formulation such as:
\begin{quote}
``In the present formalization, the geometric projection is fully threaded
through to the period theorems on both the multiset and set sides.''
\end{quote}

\item \textbf{The irrational-case unboundedness argument could be sharpened.}

The new paragraph around
\[
\texttt{LaTeX/rational\_cut\_and\_project\_gap\_periods.tex:1233--1246}
\]
now explicitly addresses unboundedness, which is a good improvement.  For
maximum precision, it would be better to state why the accepted points are
unbounded in both the positive and negative physical directions.

One concise way to do this is to use the identity
\[
  \tilde p(x,y) = x + a y
  = (1+a^2)x + a\,s(x,y),
  \qquad s(x,y)=y-a x.
\]
Since \(s(x,y)\in W\) is bounded on accepted points, \(\tilde p(x,y)\) tends
to \(+\infty\) or \(-\infty\) whenever accepted points exist with
\(x\to+\infty\) or \(x\to-\infty\), respectively.  The needed existence
follows from the density of the positive and negative tails of the irrational
rotation modulo \(1\).

This is not a flaw in the main rational proof, but it would make the
complementary irrational proposition more robust.

\item \textbf{The Lean build succeeds, but still emits style warnings.}

Running
\[
\texttt{lake build}
\]
inside
\[
\texttt{Lean/CutAndProject}
\]
completed successfully.  The build still emits style warnings, mostly long
lines and requests for comments explaining modified \texttt{maxHeartbeats}
settings.

These warnings do not affect correctness, but a warning-light or
warning-free build would make the companion formalization feel more polished.

\end{enumerate}

\section*{Checks Performed}

\begin{enumerate}[leftmargin=*]

\item \textbf{Paper build.}

The command
\begin{quote}
\texttt{latexmk -g -pdf -interaction=nonstopmode -halt-on-error}\\
\texttt{rational\_cut\_and\_project\_gap\_periods.tex}
\end{quote}
completed successfully.  I did not see unresolved references, citation
errors, overfull boxes, or remaining \texttt{hyperref} PDF-string warnings in
the searched log output.

\item \textbf{CSV verification.}

Running
\[
\texttt{python3 src/python/verify\_paper.py}
\]
on the three cited CSV files reported:
\[
\begin{array}{rl}
\text{total rows:} & 51{,}729,\\
\text{multiset mismatches:} & 0,\\
\text{set mismatches:} & 0,\\
\text{negative set sentinels:} & 0.
\end{array}
\]
Thus the empirical verification statement is now consistent with the
repository data.

\item \textbf{Lean build.}

The command
\[
\texttt{lake build}
\]
inside \texttt{Lean/CutAndProject} completed successfully.  The current Lean
file contains the headline theorems
\[
\texttt{main\_theorem\_geometric\_concrete}
\quad\text{and}\quad
\texttt{set\_main\_theorem\_geometric\_concrete},
\]
and no occurrences of \texttt{sorry}, \texttt{admit}, or \texttt{axiom} were
found in the checked source.

\end{enumerate}

\section*{Overall Assessment}

The rational core of the paper now looks mathematically solid to me.  The
main reduction
\[
  (x,y)\mapsto (r,\tilde p),
  \qquad
  r=\alpha x-\beta y,\quad
  \tilde p=\beta x+\alpha y,
\]
the residue formula
\[
  c_r \equiv -\alpha\beta^{-1}r \pmod D,
\]
and the period-minimality argument via the stabilizer of a proper cyclic
interval fit together coherently.

I would treat the remaining findings as arXiv polish rather than mathematical
blockers.  The most important small cleanup is the stale commit reference in
the Lean correspondence table, because it affects reproducibility.

\end{document}
